\documentclass[
    compress,
    10pt,
    c,
    aspectratio=1610
]{beamer}
\usefonttheme{professionalfonts}

%%%%%%%%%%%%%%%%
%  Standalone  %
%%%%%%%%%%%%%%%%
% We load currfile explicitly and very early to avoid empty stack warnings.
\usepackage{currfile}
\usepackage{standalone}
\usepackage{xstring} % standalone path magic

%%%%%%%%%%%%%%
%  Language  %
%%%%%%%%%%%%%%
\usepackage[english]{babel}
\uselanguage{English}
\languagepath{English}

%%%%%%%%%%%
%  Fonts  %
%%%%%%%%%%%
\usepackage[T1]{fontenc}
\usepackage[utf8]{inputenc}

% You might want to change this
% See http://www.tug.dk/FontCatalogue
\usepackage{arev}
\usepackage{arevmath}
\usepackage{microtype}

%%%%%%%%%%%%%%%%%%%%%%%%%%%%%
%  Used by beamer or theme  %
%%%%%%%%%%%%%%%%%%%%%%%%%%%%%
\usepackage{url}
\usepackage{listings}
\usepackage{xcolor}
\usepackage{xspace}

%%%%%%%%%%
%  Math  %
%%%%%%%%%%
\usepackage{amssymb}
\usepackage{amsfonts}
\usepackage{amsmath}
\usepackage{mathtools}
\usepackage{xfrac}
\usepackage{bm}
\usepackage{calc} % widthof

%%%%%%%%%%%%
%  Tables  %
%%%%%%%%%%%%
\usepackage{tabularx}
\usepackage{booktabs}

%%%%%%%%%%
%  Misc  %
%%%%%%%%%%
\usepackage{csquotes}
\usepackage{siunitx}

%%%%%%%%%%%%%%%%%%
%  Beamer Theme  %
%%%%%%%%%%%%%%%%%%
\usepackage{beamerthemeMA}
\setbeamercovered{invisible}
% Add slide-circles without subsections
\AtBeginSection[]{\subsection{}}

\usepackage{tikz}
\usepackage{pgfplots}
\pgfplotsset{compat=1.12}

\usetikzlibrary{calc}
\usetikzlibrary{shapes.multipart}
\usetikzlibrary{positioning}

\tikzstyle{edge} = [draw, thick, shorten >=2pt, shorten <=2pt]
\tikzstyle{directed} = [edge, arrows={-latex}, shorten >=2pt, shorten <=0pt]
\tikzstyle{marked} = [tumred]
\tikzstyle{marked edge} = [very thick, marked]
\tikzstyle{marked node} = [very thick, draw=tumred, fill=tumred!15]
\tikzstyle{every text node part} = [align=center, execute at begin node=\setlength{\baselineskip}{2ex}]

% Base Colors Corporate Design
\definecolor{tumblue}{HTML}{0065BD}
\definecolor{tumgreen}{HTML}{A2AD00}
\definecolor{tumorange}{HTML}{E37222}
\definecolor{tumivory}{HTML}{DAD7CB}
\definecolor{tumred}{HTML}{E53418} % not in Styleguide
\definecolor{tumviolet}{HTML}{69085A} % not in Styleguide

% Derived Colors
% https://kuler.adobe.com/create/color-wheel/
% https://portal.mytum.de/corporatedesign/print/styleguide/
% Grays - TUM
\definecolor{tumgray0}{HTML}{000000}
\definecolor{tumgray1}{HTML}{58585A}
\definecolor{tumgray2}{HTML}{9C9D9F}
\definecolor{tumgray3}{HTML}{D9DADB}
\definecolor{tumgray4}{HTML}{FFFFFF}

% Blues - TUM
\definecolor{tumblue0}{HTML}{003359}
\definecolor{tumblue1}{HTML}{005293}
\definecolor{tumblue2}{HTML}{0073CF}
\definecolor{tumblue3}{HTML}{64A0C8}
\definecolor{tumblue4}{HTML}{98C6EA}

% Greens - Adobe
\definecolor{tumgreen0}{HTML}{EAF900}
\definecolor{tumgreen1}{HTML}{AEBA00}
\definecolor{tumgreen2}{HTML}{8A9300}
\definecolor{tumgreen3}{HTML}{525800}

% Reds - Adobe
\definecolor{tumred0}{HTML}{F23719}
\definecolor{tumred1}{HTML}{CB2E15}
\definecolor{tumred2}{HTML}{90210F}
\definecolor{tumred3}{HTML}{65170B}

% Oranges - Adobe
\definecolor{tumorange0}{HTML}{F07824}
\definecolor{tumorange1}{HTML}{C9651E}
\definecolor{tumorange2}{HTML}{8E4715}
\definecolor{tumorange3}{HTML}{63320F}

% See http://tex.stackexchange.com/a/51638

\newcounter{jumping}
\resetcounteronoverlays{jumping}

\makeatletter
\tikzset{
  stop jumping/.style={
    execute at end picture={%
      \stepcounter{jumping}%
      \immediate\write\pgfutil@auxout{%
        \noexpand\jump@setbb{\the\value{jumping}}{\noexpand\pgfpoint{\the\pgf@picminx}{\the\pgf@picminy}}{\noexpand\pgfpoint{\the\pgf@picmaxx}{\the\pgf@picmaxy}}
      },
      \csname jump@\the\value{jumping}@maxbb\endcsname
      \path (\the\pgf@x,\the\pgf@y);
      \csname jump@\the\value{jumping}@minbb\endcsname
      \path (\the\pgf@x,\the\pgf@y);
    },
  }
}
\def\jump@setbb#1#2#3{%
  \@ifundefined{jump@#1@maxbb}{%
    \expandafter\gdef\csname jump@#1@maxbb\endcsname{#3}%
  }{%
    \csname jump@#1@maxbb\endcsname
    \pgf@xa=\pgf@x
    \pgf@ya=\pgf@y
    #3
    \pgfmathsetlength\pgf@x{max(\pgf@x,\pgf@xa)}%
    \pgfmathsetlength\pgf@y{max(\pgf@y,\pgf@ya)}%
    \expandafter\xdef\csname jump@#1@maxbb\endcsname{\noexpand\pgfpoint{\the\pgf@x}{\the\pgf@y}}%
  }
  \@ifundefined{jump@#1@minbb}{%
    \expandafter\gdef\csname jump@#1@minbb\endcsname{#2}%
  }{%
    \csname jump@#1@minbb\endcsname
    \pgf@xa=\pgf@x
    \pgf@ya=\pgf@y
    #2
    \pgfmathsetlength\pgf@x{min(\pgf@x,\pgf@xa)}%
    \pgfmathsetlength\pgf@y{min(\pgf@y,\pgf@ya)}%
    \expandafter\xdef\csname jump@#1@minbb\endcsname{\noexpand\pgfpoint{\the\pgf@x}{\the\pgf@y}}%
  }
}
\makeatother

% Sets
\newcommand{\N}{\mathbb{N}}
\newcommand{\Z}{\mathbb{Z}}
\newcommand{\R}{\mathbb{R}}
\newcommand{\Es}{\mathcal{S}}
\newcommand{\Ah}{\mathcal{A}}
\newcommand{\Tee}{\mathcal{T}}
\newcommand{\D}{\mathcal{D}}

% Constants
\newcommand{\CM}{\text{CM}}

% Random Variable
\newcommand{\rv}[1]{\bm{#1}}
\newcommand{\map}[1]{#1^{\text{MAP}}}
% Matrix
\newcommand{\mat}[1]{\bm{#1}}
\newcommand{\Eye}{\mathrm{I}}
% Pseudo Inputs
\newcommand{\ps}[1]{\bar{#1}}
\newcommand{\psmat}[1]{\ps{\mat{#1}}}

% Nicer empty set
\renewcommand{\emptyset}{\varnothing}

% Math operators
% General
\DeclareMathOperator*{\argmax}{argmax}
\DeclareMathOperator*{\argmin}{argmin}
\DeclareMathOperator{\atanTwo}{atan2}
\DeclareMathOperator{\sgn}{sgn}
\DeclareMathOperator{\diag}{diag}
\DeclareMathOperator*{\maximize}{maximize}
\DeclareMathOperator*{\minimize}{minimize}
\DeclareMathOperator{\subjectto}{subject\ to}
\DeclareMathOperator{\Oh}{\mathcal{O}}
\newcommand{\Powerset}[1]{2^{#1}}
\newcommand*{\diff}{\mathop{}\!\mathrm{d}}
\DeclarePairedDelimiter{\abs}{\vert}{\vert}

% Bicycle
\DeclareMathOperator{\DynBicycle}{\Dyn_{\text{bicycle}}}
\DeclareMathOperator{\RwdBicycle}{\Rwd_{\text{bicycle}}}
\DeclareMathOperator{\RwdBicycleProb}{\hat{\Rwd}_{\text{bicycle}}}
\DeclareMathOperator{\Goal}{\Tee_{\text{goal}}}
\DeclareMathOperator{\Fallen}{\Tee_{\text{fallen}}}
\DeclareMathOperator{\squeeze}{\rho}
\newcommand{\Trajectory}{\mat{\tau}}

% RL
\DeclareMathOperator{\Dists}{\mathcal{P}}
\DeclareMathOperator{\Rwd}{\mathbf{R}}
\DeclareMathOperator{\Vlu}{\mathbf{V}}
\DeclareMathOperator{\Dyn}{\mathbf{f}}
\DeclareMathOperator{\AVlu}{\hat{\Vlu}}

% PSO
\DeclareMathOperator{\Pe}{\mathcal{P}}
\DeclareMathOperator{\Pos}{\mathbf{x}}
\DeclareMathOperator{\Vel}{\mathbf{v}}
\DeclareMathOperator{\Neigh}{\mathcal{N}}
\DeclareMathOperator{\Cog}{\mathbf{y}}
\DeclareMathOperator{\Soc}{\hat{\Cog}}

% Probabilities
\DeclareMathOperator{\E}{\mathbb{E}}
\DeclareMathOperator{\cov}{cov}
\DeclareMathOperator{\var}{var}
\DeclareMathOperator{\corr}{\varrho}
\DeclareMathOperator{\p}{p}
\DeclareMathOperator{\K}{\mathcal{K}}
\DeclareMathOperator{\Q}{\mathcal{Q}}
\DeclareMathOperator{\Norm}{\mathcal{N}}
\DeclareMathOperator{\Uni}{\mathbb{U}}
\DeclareMathOperator{\Ind}{\mathbb{I}}
\DeclareMathOperator{\GP}{\mathcal{GP}}
\DeclareMathOperator{\SPGP}{\mathcal{SPGP}}
\DeclareMathOperator{\Ell}{\mathcal{L}}
\DeclareMathOperator{\X}{\mathcal{X}}

\providecommand\given{}
\DeclarePairedDelimiterX{\Cond}[1]{(}{)}{
\renewcommand\given{%
  \nonscript\:
  \delimsize\vert
  \nonscript\:
  \mathopen{}
  \allowbreak}
#1
}
\newcommand{\Prob}{\p\Cond}
\newcommand{\Gaussian}{\Norm\Cond}
\newcommand{\Uniform}{\Uni\Cond}
\newcommand{\Indicator}{\Ind\Cond}

\DeclarePairedDelimiterXPP{\Moment}[2]{#1}{[}{]}{}{
\renewcommand\given{%
  \nonscript\:
  \delimsize\vert
  \nonscript\:
  \mathopen{}
  \allowbreak}
#2
}

\providecommand\with{}
\DeclarePairedDelimiterX{\Set}[1]{\{}{\}}{
\renewcommand\with{%
  \nonscript\:
  \delimsize\vert
  \nonscript\:
  \mathopen{}
  \allowbreak}
#1
}

\DeclarePairedDelimiter{\Bicycle}{(}{)}


\title{Incorporating Uncertainty into Reinforcement Learning through Gaussian Processes}
\author{Markus Kaiser}
\subtitle{Master's Thesis}
\author{\href{mailto:markus.kaiser@in.tum.de}{Markus Kaiser}}
\titlegraphic{\parbox{.95\textwidth}{%
    \raisebox{-.5\height}{\includegraphics[height=1.5cm,width=.33\linewidth,keepaspectratio]{tum}}%
    \hfill%
    \raisebox{-.5\height}{\includegraphics[height=1.5cm,width=.33\linewidth,keepaspectratio]{kth}}%
    \hfill%
    \raisebox{-.5\height}{\includegraphics[height=1.5cm,width=.33\linewidth,keepaspectratio]{siemens}}%
}}
\date{4. July 2016}

% \includeonlyframes{solution_psop,solution_reward}

\begin{document}

\begin{frame}[plain, noframenumbering]
    \titlepage
\end{frame}

\begin{frame}[noframenumbering]
    \frametitle{Outline}

    \tableofcontents
\end{frame}

\section{The Bicycle Benchmark}
\begin{frame}[label=bicycle_drive]
    \frametitle{Bicycle Benchmark -- Driving}

    \begin{columns}
        \begin{column}[c]{0.5\textwidth}
            \centering
            \includestandalone{figures/bicycle_from_above}
        \end{column}
        \begin{column}[c]{0.5\textwidth}
            \centering
            \vspace{1cm}
            \uncover<6-> {
                \includestandalonewithpath{figures/bicycle_drive_goal}
            }
        \end{column}
    \end{columns}
\end{frame}

\begin{frame}[label=bicycle_balance]
    \frametitle{Bicycle Benchmark -- Balancing}

    \vfill
    \begin{columns}
        \begin{column}{0.45\textwidth}
            \centering
            \includestandalone{figures/bicycle_from_behind}
        \end{column}
        \begin{column}{0.55\textwidth}
            \centering
            \uncover<4-> {
                \includestandalonewithpath{figures/bicycle_drive_balance}
            }
        \end{column}
    \end{columns}
\end{frame}

\section{Gaussian Processes}
\begin{frame}[label=gp_non_parametric_models]
    \frametitle{Bayesian Non-Parametric Models}

    \begin{overlayarea}{\textwidth}{8\baselineskip}
        \only<1> {
            \begin{block}{Data}
                \begin{itemize}
                    \item Observe a \structure{noisy} data set $\D = \Set{\left( \mat{x_i}, \mat{y_i} \right)}$
                    \item Assume that $\mat{y_i} = f(\mat{x_i}) + \Gaussian{0, \sigma_n^2}$
                \end{itemize}
            \end{block}
        }
        \only<2> {
            \begin{block}{Non-Bayesian Parametric Models}
                \begin{itemize}
                    \item Assume \structure{structure} about $f$, such as $f(\mat{x}) = \mat{W}\mat{x}$
                    \item Find a \structure{single} $\mat{W}$ which is optimal w.r.t.\ some criterion
                    \item Everything is \structure{deterministic}
                \end{itemize}
            \end{block}
        }
        \only<3-4> {
            \begin{block}{Bayesian Parametric Models}
                \begin{itemize}
                    \item Assume \structure{structure} about $f$, such as $f(\mat{x}) = \mat{W}\mat{x}$
                    \item Find a \structure{distribution} $\Prob{\mat{W} \given \D}$ of plausible parameters
                    \item We can consider all parameters at once
                \end{itemize}
            \end{block}
        }
        \only<5-> {
            \begin{block}{Bayesian Non-Parametric Models}
                \begin{itemize}
                    \item $f$ can be an \structure{arbitrary function}
                    \item Find a \structure{distribution} $\Prob{f \given \D}$ over functions
                    \item We can consider all \structure{functions} at once
                \end{itemize}
            \end{block}
        }
    \end{overlayarea}
    \vfill
    \centering
    \includestandalonewithpath{figures/non_parametric_models}
\end{frame}

\begin{frame}[label=gp_definition]
    \frametitle{Gaussian Processes}

    \begin{definition}[Gaussian Process]
        A \structure{Gaussian Process (GP)} is a collection of random variables $\Set{\rv{F_x} \with x \in X}$, any finite subset of which has a joint Gaussian distribution.
    \end{definition}
    \begin{itemize}
        \item Extension of Gaussians to (infinite) function spaces
        \item $\rv{F_x}$ models the function value $f(x)$ of some $f : X \to Y$
    \end{itemize}
    \vfill
    \pause
    \begin{block}{Mean and Kernel Functions}
        A GP is \structure{completely determined} by
        \begin{description}[\quad Kernel function]
            \item[Mean function] \makebox[\widthof{$\K(x, x^\prime)$}][r]{$\mu_f(x)$}${} = \Moment{\E}{f(x)}$
            \item[Kernel function] $\K(x, x^\prime) = \Moment{\cov}{f(x), f(x^\prime)}$
        \end{description}
    \end{block}
    \begin{itemize}
        \item The kernel encodes the \structure{prior assumptions} about the function
    \end{itemize}
\end{frame}

\begin{frame}[label=gp_hyperparameters]
    \frametitle{Posterior and Hyperparameters}

    \centering
    \includestandalonewithpath{figures/gp_hyperparameters}
\end{frame}

\section{Uncertainties in Reinforcement Learning}
\begin{frame}[label=solution_architecture]
    \frametitle{Solution Architecture}

    \centering
    \includestandalone{figures/agent_environment}
\end{frame}

\begin{frame}[label=solution_psop]
    \frametitle{PSO-P}

    \begin{definition}[PSO-P]
        \label{def:psop}
        The \structure{Particle Swarm Optimization-Policy (PSO-P)} chooses actions via optimization of the expected accumulated reward.
        \begin{align}
            \pi_{\text{PSO-P}}(\mat{s}) &\coloneqq \mat{a^\ast}_0\text{, where} \\
            \mat{a^\ast} &\in \argmax_{\mat{a} \in \Ah^T} \left\{ \Moment*{\E}{\sum_{t=1}^T \gamma^t \Rwd(\mat{s_t}) \given \GP, \mat{s_0} = \mat{s}, \mat{a_0}, \dots, \mat{a_{T-1}}} \right\}
        \end{align}
    \end{definition}
    \begin{itemize}
        \item PSO is a gradient-free heuristic
        \item Directly exploits the transition models
    \end{itemize}
\end{frame}

\begin{frame}[label=solution_reward]
    \frametitle{Reward Function}

    \begin{definition}[Bicycle Reward Function]
        The \structure{bicycle reward function} is defined as
        \begin{align}
            \RwdBicycle(\mat{s}) &\coloneqq \begin{cases}
                    2 & \text{if goal reached} \\
                    0 & \text{if fallen down} \\
                    c\cdot\Gaussian{\Delta^\psi_{\mat{s}} \given 0, \sigma_{\text{angle}}^2} & \text{otherwise}
                \end{cases}
        \end{align}
    where $\Delta^\psi_{\mat{s}}$ denotes the \structure{angle towards the goal}.
    \end{definition}
    \vspace{2ex}
    \centering
    \includestandalone{figures/saturating_reward}
\end{frame}

\begin{frame}[label=solution_map]
    \frametitle{MAP without Uncertainties}

    \centering
    \includestandalonewithpath{figures/map}
\end{frame}

\begin{frame}[label=solution_map_predictions]
    \frametitle{MAP Predictions}

    \centering
    \includestandalonewithpath{figures/map_predictions}
\end{frame}

\begin{frame}[label=solution_os]
    \frametitle{One-Step Uncertainties}

    \centering
    \includestandalonewithpath{figures/one_step}
\end{frame}

\begin{frame}[label=solution_os_predictions]
    \frametitle{One-Step Predictions}

    \centering
    \includestandalonewithpath{figures/one_step_uncertainties}
\end{frame}

\begin{frame}[label=solution_ms]
    \frametitle{Multi-Step Uncertainties}

    \centering
    \includestandalonewithpath{figures/linearization}
\end{frame}

\begin{frame}[label=solution_ms_predictions]
    \frametitle{Multi-Step Predictions}

    \centering
    \includestandalonewithpath{figures/linearization_predictions}
\end{frame}

\section{Results}
\begin{frame}[label=results_qualitative]
    \frametitle{Qualitative Results}

    \centering
    \includestandalonewithpath{figures/xy_trajectories}
\end{frame}

\begin{frame}[label=results_quantitative]
    \frametitle{Quantitative Results}

    \centering
    \vfill
    \begin{columns}
        \begin{column}[c]{0.5\textwidth}
            \centering
            \uncover<3> {
                \includestandalonewithpath{figures/goal_percentage}
            }
        \end{column}
        \begin{column}[c]{0.5\textwidth}
            \centering
            \begin{tabularx}{\textwidth}{X>{\centering\hsize=.4\hsize}X>{\centering\hsize=.4\hsize}X>{\centering\hsize=.4\hsize\arraybackslash}X}
                \toprule
                Metric & MAP & OS & MS \\
                \midrule
                Trajectories & 9660 & 9660 & 9660
                \uncover<2-> {
                    \\
                    \addlinespace
                    Success Rate & \SI[mode=text]{53.4}{\percent} & \SI[mode=text]{63.8}{\percent} & \textbf{\SI[mode=text,detect-weight]{75.8}{\percent}} \\
                    \addlinespace
                    Time to Goal \\
                    Mean & \textbf{59.9} & 62.0 & 66.5 \\
                    Median & \textbf{60} & \textbf{60} & 63
                }
                \\
                \bottomrule
            \end{tabularx}
        \end{column}
    \end{columns}
\end{frame}

\appendix
\section{Additional Material}
\begin{frame}[noframenumbering]
    \vfill
    \centering
    \begin{beamercolorbox}[sep=8pt,center]{title}
        \usebeamerfont{title}Additional Material\par%
    \end{beamercolorbox}
    \vfill
\end{frame}

\begin{frame}[label=gp_sparse, noframenumbering]
    \frametitle{Sparse Gaussian Processes}

    \centering
    \includestandalonewithpath{figures/spgp}
\end{frame}

\end{document}
